Fix \(c\), \(k\), \(\ell\), and \(z\in\mathcal S_k\).  By \(\pi_\ell\in(0,1)\), the Bernoulli-unit randomization, and independence of labels from schedules,
\[
 \Pr(L_c=\ell,Z_c=z)
 =p_\ell\pi_\ell^{m_k}(1-\pi_\ell)^{n-m_k}
 =\frac{p_\ell B_{k\ell}}{M_k}.
\]
Summing this equality over \(\ell\in\mathcal K\), and using \(q=Bp\), gives
\[
 \Pr(Z_c=z)=\frac{1}{M_k}\sum_{\ell=1}^Kp_\ell B_{k\ell}=\frac{q_k}{M_k}.
\]
Every entry of \(B\) is positive because \(1\leq m_k\leq n-1\) and \(0<\pi_\ell<1\); hence \(q_k>0\), even when \(p\) is on the boundary.  Division by \(q_k/M_k\) is therefore valid and Bayes' formula yields
\[
 \Pr(L_c=\ell\mid Z_c=z)=\frac{p_\ell B_{k\ell}}{q_k}.
\]
The value \(q_k/M_k\) does not depend on \(z\in\mathcal S_k\), so summing it over the \(M_k\) cells shows both \(\Pr(Z_c\in\mathcal S_k)=q_k\) and uniformity conditional on that event.  The slices are disjoint because the counts \(m_k\) are distinct.  Finally, the cluster records are independent under the iid label and schedule assumptions and the within-cluster randomization.  Thus each cluster falls in exactly one of the \(K+1\) categories \(\mathcal S_1,\ldots,\mathcal S_K,(\bigcup_k\mathcal S_k)^c\), with probabilities \(q_1,\ldots,q_K,1-\sum_kq_k\), proving the multinomial assertion.